% SSSUP-FONTE-STORICA: fonti-locali/originali/2014-2020.zip
% Anno 2014, esercizio 2. Trasposizione dal sorgente LaTeX storico.
% Verificare editorialmente la trasposizione prima di considerarla definitiva.
\renewcommand{\Anno}{2014}
\renewcommand{\NumeroProblema}{2}

\renewcommand{\Testo}{
Nella regione di Toscolandia esistono tre città: FI, PI, LI, disposte ai vertici di un triangolo rettangolo avente i cateti di lunghezza 80 km (tra PI e FI) e 20 km (tra PI e LI). Le città sono tutte collegate tra loro da ottime strade rettilinee, ma esiste una sola ferrovia, che collega PI a FI, senza stazioni intermedie. I cittadini di PI e di LI devono recarsi regolarmente a FI, capoluogo della regione, e possono scegliere se andare in auto o in treno; i cittadini di LI, se vogliono, possono andare a PI in auto e prendere poi a PI il treno, oppure andare direttamente in auto. Si tenga conto che:
\begin{itemize}
\item  il numero di cittadini che si recano a FI da PI e da LI è lo stesso;
\item si possono trascurare gli altri tipi di viaggiatori (ad es. cittadini di FI o turisti);
\item ogni cittadino, a parità di costo totale individuale, per comodità sceglie il treno (o auto + treno nel caso di LI);
\item i viaggi in auto sono effettuati con il solo guidatore, senza passeggeri;
\item il costo per km di un tragitto in auto è di $p$ Euro ($p=0.15$ per una simulazione realistica).
\end{itemize}

\begin{enumerate}
\item[(a)] La compagnia ferroviaria regionale vuole decidere il prezzo del biglietto PI-FI in modo da massimizzare le entrate; quale prezzo deciderà?
\item[(b)] Con gli stessi dati di sopra determinare il prezzo che massimizza le entrate nel caso in cui invece che 20 km la distanza PI-LI sia $d$.
\end{enumerate}
}

\renewcommand{\Risposta}{
% SSSUP-BOZZA-AI: Risposta; autore=ChatGPT; data=2026-11-10; stato=DA-VERIFICARE
% SSSUP-PROVENIENZA: ChatGPT (OpenAI), bozza AI 2014 del 2026-11-10; ricostruzione indipendente dal Testo storico; revisione umana pendente.
\begin{enumerate}
\item[(a)] $p(\sqrt{6800}-20)$ euro, ossia circa $9.3693$ euro
se $p=0.15$.
\item[(b)] Posto $b(d)=p(\sqrt{6400+d^2}-d)$, il prezzo
ottimale è $b(d)$ per $0\leq d<60$, oppure $80p$ per
$d>60$. Per $d=60$ entrambi i prezzi sono ottimali.
\end{enumerate}
}
\renewcommand{\Aiutino}{
% SSSUP-BOZZA-AI: Aiutino; autore=ChatGPT; data=2026-11-10; stato=DA-VERIFICARE
% SSSUP-PROVENIENZA: ChatGPT (OpenAI), bozza AI 2014 del 2026-11-10; ricostruzione indipendente dal Testo storico; revisione umana pendente.
Confrontare, separatamente per gli abitanti di PI e di LI,
il costo del treno con quello dell'auto. Si ottengono due
soglie per il prezzo del biglietto: il ricavo cambia formula
quando si supera una di esse. Basta confrontare i ricavi
massimi ottenibili nelle due fasce di prezzo.
}
\renewcommand{\SoluzioneUfficiale}{}
\renewcommand{\SoluzioneRedazionale}{
% SSSUP-BOZZA-AI: SoluzioneRedazionale; autore=ChatGPT; data=2026-11-10; stato=DA-VERIFICARE
% SSSUP-PROVENIENZA: ChatGPT (OpenAI), bozza AI 2014 del 2026-11-10; ricostruzione indipendente dal Testo storico; revisione umana pendente.
Sia $x$ il prezzo del biglietto PI--FI, e sia $N$ il numero
di viaggiatori provenienti da ciascuna delle due città.

Per un abitante di PI, il viaggio in auto costa $80p$ euro.
Poiché in caso di parità si preferisce il treno, tutti gli
abitanti di PI scelgono il treno se e solo se
\[
 x\leq A:=80p.
\]

Per un abitante di LI il tragitto diretto in auto misura
$\sqrt{80^2+d^2}=\sqrt{6400+d^2}$ chilometri, mentre
l'alternativa auto fino a PI e treno fino a FI costa
$dp+x$. Pertanto i viaggiatori di LI prendono il treno
se e solo se
\[
 x\leq B:=p(\sqrt{6400+d^2}-d).
\]
Per $d\geq0$ si ha $0<B\leq A$.

Il ricavo ferroviario è dunque
\[
 R(x)=
 \begin{cases}
  2Nx,&0\leq x\leq B,\\
  Nx,&B<x\leq A,\\
  0,&x>A.
 \end{cases}
\]
(In caso di coincidenza delle soglie, la fascia centrale
è vuota.) Le uniche possibilità per il massimo sono
$R(B)=2NB$ e $R(A)=NA$. Conviene scegliere $B$ quando
$2B\geq A$, cioè quando
\[
 2(\sqrt{6400+d^2}-d)\geq80
 \quad\Longleftrightarrow\quad d\leq60.
\]
Per $d<60$ il prezzo ottimale è quindi $B$; per $d>60$
è $A$; per $d=60$ entrambi danno lo stesso ricavo massimo.

Nel caso (a), $d=20<60$ e dunque
\[
 x=p(\sqrt{6800}-20).
\]
Con $p=0.15$ il risultato è circa $9.3693$ euro.
Se il prezzo deve essere espresso in centesimi, conviene
arrotondare per difetto a $9.36$ euro: arrotondando per
eccesso a $9.37$ euro i viaggiatori di LI non prenderebbero
il treno.
}
\renewcommand{\SoluzioniAlternative}{}
\renewcommand{\Osservazioni}{}
\renewcommand{\Commenti}{}
